\bookchapter{Глава 21: Глубокое погружение в мобильные сети и информационную безопасность}

\chapterintro{Каждый день мы пользуемся мобильным интернетом, оплачиваем покупки, общаемся в мессенджерах — всё так удобно! Но задумывались ли вы, как часто ваши переписки, платёжные пароли, фотографии… «голышом» гуляют по сети? Честно говоря, большинство людей вообще об этом не думает, но когда случается беда — например, угон аккаунта или списание с банковской карты — тогда уже поздно жалеть. В этой главе мы подробно разберём, как сделать свою цифровую жизнь менее «прозрачной».}

\booksection{VPN: «мантия-невидимка» в сети}

VPN — сокращение от Virtual Private Network (виртуальная частная сеть). Представьте, что между вашим домом и нужным сайтом проложен эксклюзивный зашифрованный подземный тоннель. Вы идёте по тоннелю, а посторонние не видят, что у вас в руках и куда вы направляетесь.

\begin{itemize}
  \item \textbf{Обычный интернет}: ваш телефон общедоступная сеть целевой сайт. На любом промежуточном маршрутизаторе или коммутаторе кто-то может «подглядывать» за отправляемыми данными. Особенно в кафе или аэропорту через публичный Wi-Fi — грубо говоря, вы «голышом».
  \item \textbf{С VPN}: ваш телефон зашифрованный тоннель VPN сервер VPN целевой сайт. Данные зашифрованы и упакованы, «подсматривающий» видит лишь бессмысленный набор символов. Кроме того, целевой сайт видит IP-адрес сервера VPN, а не ваш реальный IP.
\end{itemize}

Что же реально даёт VPN? Если коротко: во-первых,\textbf{защищает вашу приватность на публичном Wi-Fi}; во-вторых,\textbf{обходит географические ограничения}(например, контент, доступный только в определённых странах); в-третьих,\textbf{не даёт вашему провайдеру (ISP) отслеживать ваши привычки просмотра}. Но предупреждаю:\textbf{не пользуйтесь бесплатными VPN!}Бесплатное часто означает, что ваши данные — их «товар». Один мой знакомый пользовался бесплатным VPN, и телефон начал сам по себе показывать рекламу — оказалось, VPN продавал его данные рекламодателям. Так что выбирайте либо платные VPN крупных брендов (NordVPN, ExpressVPN), либо не пользуйтесь вовсе.

И ещё: VPN — не панацея. Он защищает только данные «в пути». Если целевой сайт сам небезопасен или на вашем телефоне уже сидит вредоносное ПО, VPN не поможет.

\booksection{«Сладкая ловушка» публичного Wi-Fi — опаснее, чем вы думаете}

Может показаться: ну что там, подключился к бесплатному Wi-Fi, какой риск? Расскажу реальный случай: исследователи безопасности поставили эксперимент — развернули в аэропорту точку доступа под названием «Starbucks\_Free», и менее чем за час к ней подключились более ста человек. На самом деле эту точку создали сами исследователи, и весь\textbf{нешифрованный трафик}, проходящий через неё — например, содержимое HTTP-страниц, коммуникация приложений без TLS — они реально могли перехватывать и просматривать.

Но важно понимать ключевое условие: сегодня WeChat, мобильный банкинг, Alipay и большинство привычных вам приложений уже используют шифрование HTTPS. Проще говоря, даже если злоумышленник перехватит пакет данных, он увидит лишь бессмысленный набор символов,\textbf{и не сможет напрямую прочитать содержимое переписки или платёжный пароль}. Однако он всё равно узнает, какие сайты (домены) вы посещали. Если какое-то устаревшее приложение или сайт ещё использует открытую HTTP-передачу, содержимое будет полностью раскрыто. Такой тип атаки называется «атака человек посередине» (Man-in-the-Middle): злоумышленник «вклинивается» между вами и настоящим Wi-Fi-маршрутизатором. Ещё страшнее то, что у многих установлен «автоматическое подключение к известным Wi-Fi» — злоумышленнику достаточно назвать свою точку так же, как сеть, к которой вы раньше подключались (например, «ChinaNet» или «CMCC»), и телефон сам к ней подключится.

Как защититься? Вот несколько практических советов:

\begin{itemize}
  \item \textbf{Смотрите на HTTPS}: адрес начинается с https:// — лишняя «s» означает Secure. На любом сайте, куда нужно вводить пароль, если видите http:// (без s) — немедленно остановитесь и ничего не вводите! Современные браузеры помечают не-HTTPS сайты предупреждением «небезопасно» — увидели такой значок, будьте настороже.
  \item \textbf{Отключите «автоматическое подключение»}: и на iPhone, и на Android в настройках Wi-Fi выключите «автоподключение» или «автовступление». Лучше каждый раз нажимать вручную, чем позволять телефону решать за вас.
  \item \textbf{Чувствительные операции — только через мобильный трафик}: переводы, вход в мобильный банк, ввод номера удостоверения… эти операции выполняйте честно через свой 4G/5G, не экономьте на трафике.
  \item \textbf{«Забывайте» редко используемые Wi-Fi}: вернулись из командировки или отпуска — «забудьте» временные сети отеля и аэропорта. Тогда телефон не подключится автоматически к подставной точке с таким же именем.
\end{itemize}

\booksection{Как HTTPS вас защищает? — SSL/TLS-рукопожатие}

Мы упомянули HTTPS, но как именно он работает? Каждый раз при посещении HTTPS-сайта браузер и сервер выполняют «SSL/TLS-рукопожатие», и процесс примерно таков:

\begin{enumerate}
  \item Ваш браузер говорит: «Привет, сервер, я хочу безопасно общаться, вот список моих криптографических возможностей (поддерживаемые алгоритмы шифрования и т. д.).»
  \item Сервер отвечает: «Хорошо, я выбрал из твоего списка эту схему шифрования, вот мой цифровой сертификат — в нём мой открытый ключ, подписанный и сертифицированный таким-то Удостоверяющим центром (CA).»
  \item Браузер проверяет подлинность сертификата: выдан ли доверенным CA? Не истёк ли срок? Соответствует ли домен?
  \item После успешной проверки обе стороны с помощью открытого ключа согласовывают «сеансовый ключ», и дальнейший обмен шифруется этим ключом.
\end{enumerate}

Это похоже на визит в банк: вы предъявляете удостоверение (сертификат), кассир проверяет вашу личность, затем вы договариваетесь о кодовом слове, и дальнейший разговор понимаете только вы вдвоём. Даже если кто-то подслушивает — он слышит лишь зашифрованный набор символов.

Но учтите: HTTPS гарантирует лишь шифрование «в процессе передачи», но не гарантирует, что целевой сайт — «хороший». Фишинговый сайт тоже может использовать HTTPS — значок замка означает лишь, что вы подключены к «какому-то сайту», а не к тому, который вы хотели. Поэтому всегда самостоятельно проверяйте правильность адреса.

\booksection{Безопасность DNS: ваш «сетевой телефонный справочник» в безопасности?}

DNS (система доменных имён) — это телефонный справочник интернета: вы вводите «baidu.com», DNS находит соответствующий IP-адрес (например, 110.242.68.66), и только тогда браузер может подключиться. Проблема в том, что традиционные DNS-запросы\textbf{открытым текстом}и\textbf{не шифруются}! То есть ваш ISP, владелец Wi-Fi или любой промежуточный узел может видеть, какие сайты вы посещали.

Ещё хуже — «DNS-перехват»: злоумышленник подменяет DNS и направляет «baidu.com» на фишинговый сайт. Вы вводите правильный адрес, но попадаете на поддельную страницу входа — стоит ввести логин и пароль, и всё пропало.

К счастью, теперь есть решения с шифрованием DNS:

\begin{itemize}
  \item \textbf{DNS over HTTPS (DoH)}: DNS-запросы прячутся внутри HTTPS-трафика, смешиваясь с обычным трафиком просмотра — невозможно отличить, какие пакеты содержат DNS-запросы.
  \item \textbf{DNS over TLS (DoT)}: DNS-запросы защищаются протоколом TLS на отдельном порту; сетевой администратор видит, что вы используете DoT, но не видит содержимое.
\end{itemize}

Лично я рекомендую использовать в качестве DNS-сервера Cloudflare 1.1.1.1 или Google 8.8.8.8 — оба поддерживают шифрование DNS. На Android в настройках «Частный DNS» просто укажите «one.one.one.one» или «dns.google» — очень просто. На iPhone чуть сложнее: нужно установить профиль конфигурации или настроить через приложение, но это стоит сделать.

\booksection{Действительно ли «режим инкогнито» делает вас невидимым?}

Сразу вывод:\textbf{режим инкогнито — это шторы в вашем доме, а не мантия-невидимка.}

Что он делает: не оставляет следов локально — после закрытия окна инкогнито браузер не сохраняет историю просмотров, cookies, данные форм, поисковые запросы. Если вы пользуетесь чужим компьютером или покупаете подарок и не хотите, чтобы семья увидела историю — режим инкогнито уместен.

Чего он НЕ делает:

\begin{itemize}
  \item Посещённый вами сайт\textbf{всё равно видит, что вы пришли}— ваш IP-адрес, информацию об устройстве, поведение на странице сайту полностью видны.
  \item Ваш провайдер\textbf{всё равно видит, какие сайты вы посещали}, просто не видит конкретное содержимое (если используется HTTPS).
  \item Сетевой администратор компании/школы\textbf{по-прежнему может контролировать}вашу активность в сети.
  \item Скачанные файлы и добавленные закладки\textbf{сохраняются}и не исчезают при закрытии окна.
\end{itemize}

Так что запомните: режим инкогнито — «локальная стиральная машина», а не «сетевая мантия-невидимка». Если вам действительно нужно скрыть следы в сети, понадобится комбинация VPN + браузер Tor, но для повседневного использования достаточно понимать эти ограничения.

\booksection{Конфиденциальность браузера: не дайте браузеру себя продать}

Вы можете не осознавать, но сам браузер — огромный сборщик информации. Cookies, история просмотров, кэш, автозаполнение, доступ к местоположению… каждый пункт может раскрыть вашу приватность. Несколько советов:

\begin{itemize}
  \item \textbf{Регулярно очищайте данные браузера}: время от времени заходите в настройки и чистите cookies, кэш, историю. Не столько ради защиты «от слежки» — очистка cookies уже не особо помогает — сколько ради того, чтобы телефон не тормозил. Кэш браузера может занимать несколько гигабайт.
  \item \textbf{Забудьте про фантазию «Не отслеживать»}: опция «Отправлять запрос Do Not Track» — это по сути «просьба» сайту не отслеживать вас. Но послушает ли сайт — совершенно другой вопрос, большинство просто игнорирует. Не надейтесь на это.
  \item \textbf{Используйте браузеры с защитой приватности}: Firefox Focus, браузер DuckDuckGo, Brave — эти браузеры по умолчанию блокируют трекеры и рекламу, гораздо чище, чем стандартные настройки Chrome и Safari.
  \item \textbf{Поисковик тоже важен}: используйте DuckDuckGo вместо Baidu/Google — он не сохраняет историю поиска и не строит профиль пользователя. Честно говоря, результаты поиска не такие точные, как у Google, но для повседневных запросов достаточно.
  \item \textbf{Проверьте разрешения приложений}: загляните в настройки телефона и посмотрите, какие разрешения выданы браузеру. Ему действительно нужен доступ к вашему местоположению, микрофону, контактам? Скорее всего, нет.
\end{itemize}

\booksection{Практический чек-лист безопасности: начните защищать себя прямо сейчас}

Многое сказано, а теперь — простой и прямолинейный чек-лист. Потратьте полчаса, и безопасность вашего телефона качественно изменится:

\begin{enumerate}
  \item \textbf{Отключите автоподключение Wi-Fi}: на iPhone — Настройки Wi-Fi; на Android — настройки Wi-Fi.
  \item \textbf{Переключитесь на шифрованный DNS}: на Android укажите «dns.google» в качестве частного DNS; на iPhone настройте DoH 1.1.1.1 через приложение или профиль.
  \item \textbf{Включите двухфакторную аутентификацию (2FA)}: WeChat, Alipay, Apple ID, Google-аккаунт — везде, где можно, включите 2FA. SMS-коды — минимальный уровень 2FA; лучше использовать TOTP (Google Authenticator, Microsoft Authenticator).
  \item \textbf{Проверьте разрешения всех приложений}: пройдитесь по разрешениям каждого приложения, особенно: местоположение, контакты, SMS, микрофон, камера. Вы обнаружите, что приложение-фонарик запрашивает доступ к контактам — немедленно удаляйте.
  \item \textbf{Регулярно очищайте данные браузера}: поставьте напоминание в календаре — раз в месяц чистить кэш и cookies.
  \item \textbf{Не используйте кабели данных от публичных зарядных станций}: USB-порт может и заряжать, и передавать данные; вредоносная станция («Juice Jacking») может украсть данные во время зарядки. Берите с собой собственный адаптер и кабель.
\end{enumerate}

\begin{tcbpractice}{Сделайте прямо сейчас: проведите «медосмотр безопасности» вашего телефона}
\begin{enumerate}
  \item \textbf{Обновление DNS}: Android: Настройки Подключение и общий доступ Частный DNS введите «one.one.one.one» или «dns.google». iPhone: найдите в Safari «1.1.1.1 DNS profile» и установите профиль, либо скачайте приложение Cloudflare 1.1.1.1.
  \item \textbf{Аудит Wi-Fi}: откройте список Wi-Fi, «забудьте» все незнакомые, редко используемые и подозрительно названные сети.
  \item \textbf{Генеральная уборка разрешений}: зайдите в Настройки Конфиденциальность/управление разрешениями, проверьте разрешения каждого приложения. Отметьте приложения с необоснованными запросами (например, калькулятор запрашивает контакты) — рассмотрите удаление или замену.
  \item \textbf{Проверьте безопасность аккаунтов}: откройте страницу безопасности Apple ID или Google-аккаунта, проверьте наличие незнакомых устройств, подтвердите включение 2FA, проверьте доверенных контактных лиц для восстановления.
\end{enumerate}

\end{tcbpractice}
